\chapter{Matchings}
%%%%%%%%%%%%%%%%%%%%%%%%%%%%%%%%%%%%%%%%%%%%%%%%%%%%%%%%%%%%%%%%%%%%%%%%%%%%%%%
\section{Matchings}

% Generated by scripts/blueprint_restate.py from TCSlib.GraphTheory.Core.Matchings.
% Each body is the STATEMENT only, taken from the declaration's Lean docstring:
% the one-line summary plus the verbatim book statement.  Proof, proof plan,
% significance commentary and formalisation notes are excluded by construction.

\begin{definition}[A maximum matching: a matching no larger than which exists]
\lean{SimpleGraph.Subgraph.IsMaximumMatching}
\label{SimpleGraph.Subgraph.IsMaximumMatching}
\leanok
A \textbf{maximum matching}: a matching no larger than which exists.

A subset $M$ of $E$ is called a \emph{matching} in $G$ if its elements are links and no
two are adjacent in $G$; the two ends of an edge in $M$ are said to be \emph{matched
under $M$}. A matching $M$ \emph{saturates} a vertex $v$, and $v$ is said to be
\emph{$M$-saturated}, if some edge of $M$ is incident with $v$; otherwise, $v$ is
\emph{$M$-unsaturated}. If every vertex of $G$ is $M$-saturated, the matching $M$ is
\emph{perfect}. $M$ is a \emph{maximum matching} if $G$ has no matching $M'$ with $|M'|
> |M|$; clearly, every perfect matching is maximum.
\end{definition}

\begin{definition}[$M$-alternating walk: consecutive edges alternate in/out of $M$]
\lean{SimpleGraph.Walk.IsAlternatingWalk}
\label{SimpleGraph.Walk.IsAlternatingWalk}
\leanok
\texttt{M}-\textbf{alternating walk}: consecutive edges alternate in/out of \texttt{M}.

Let $M$ be a matching in $G$. An \emph{$M$-alternating path} in $G$ is a path whose
edges are alternately in $E \backslash M$ and $M$. For example, the path $v_5 v_8 v_1
v_7 v_6$ in the graph of figure 5.1$a$ is an $M$-alternating path.
\end{definition}

\begin{definition}[$M$-augmenting path: alternating, both ends $M$-unsaturated]
\lean{SimpleGraph.Walk.IsAugmenting}
\label{SimpleGraph.Walk.IsAugmenting}
\leanok
\uses{SimpleGraph.Walk.IsAlternatingWalk}
\texttt{M}-\textbf{augmenting path}: alternating, both ends \texttt{M}-unsaturated.

An \emph{$M$-augmenting path} is an $M$-alternating path whose origin and terminus
are $M$-unsaturated.
\end{definition}

\begin{definition}[A minimum covering]
\lean{SimpleGraph.IsMinimumCovering}
\label{SimpleGraph.IsMinimumCovering}
\leanok
\uses{SimpleGraph.IsCovering}
A \textbf{minimum covering}.

A covering $K$ is a \emph{minimum covering} if $G$ has no covering $K'$ with
$|K'| < |K|$ (see figure 5.4).
\end{definition}

\begin{definition}[A $k$-factor: a $k$-regular spanning subgraph]
\lean{SimpleGraph.Subgraph.IsKFactor}
\label{SimpleGraph.Subgraph.IsKFactor}
\leanok
A \texttt{k}-\textbf{factor}: a \texttt{k}-regular spanning subgraph.

\textbf{5.1.5} A $k$-\emph{factor} of $G$ is a $k$-regular spanning subgraph of $G$, and
$G$ is $k$-\emph{factorable} if there are edge-disjoint $k$-factors
$H_1, H_2, \ldots, H_n$ such that $G = H_1 \cup H_2 \cup \ldots \cup H_n$.
\end{definition}

\begin{definition}[$G$ is $k$-factorable: it decomposes into edge-disjoint $k$-factors]
\lean{SimpleGraph.IsKFactorable}
\label{SimpleGraph.IsKFactorable}
\leanok
\uses{SimpleGraph.Subgraph.IsKFactor}
\texttt{G} is \texttt{k}-\textbf{factorable}: it decomposes into edge-disjoint \texttt{k}-factors.

$G$ is $k$-\emph{factorable} if there are edge-disjoint $k$-factors
$H_1, H_2, \ldots, H_n$ such that $G = H_1 \cup H_2 \cup \ldots \cup H_n$.
\end{definition}

\begin{definition}[The $k$-cube $Q_k$]
\lean{SimpleGraph.cube}
\label{SimpleGraph.cube}
The \texttt{k}-\textbf{cube} \texttt{Q\_k}.

The $k$-cube is the graph whose vertices are the ordered $k$-tuples of $0$s and
$1$s, two vertices being joined if and only if they differ in exactly one
coordinate.

Defined in chapter 1; chapter 5 uses it only in exercise 5.1.1(a).
\end{definition}

\begin{definition}[The $m \times n$ grid graph]
\lean{SimpleGraph.gridGraph}
\label{SimpleGraph.gridGraph}
The \texttt{m $\times$ n} \textbf{grid graph}.
\end{definition}

\begin{theorem}[Berge's theorem on maximum matchings]
\lean{SimpleGraph.berge_maximum_matching}
\label{SimpleGraph.berge_maximum_matching}
\uses{SimpleGraph.Subgraph.IsMaximumMatching, SimpleGraph.Walk.IsAugmenting}
Let $G$ be a finite simple graph and let $M$ be a matching in $G$. Then $M$ is a maximum
matching if and only if $G$ contains no $M$-augmenting path.
\end{theorem}

\begin{theorem}[Hall's marriage theorem for bipartite matchings]
\lean{SimpleGraph.hall_bipartite_matching}
\label{SimpleGraph.hall_bipartite_matching}
Let $G$ be a finite simple graph whose vertex set admits a bipartition into parts $X$
and $Y$, so that every edge of $G$ joins a vertex of $X$ to a vertex of $Y$. For a set
$S$ of vertices write $N(S) = \bigcup_{v \in S} N(v)$ for the union of the neighborhoods
of the vertices in $S$. Then $G$ has a matching that saturates $X$ — a set of pairwise
disjoint edges covering every vertex of $X$ — if and only if $\abs{S} \le \abs{N(S)}$
for every $S \subseteq X$.
\end{theorem}

\begin{theorem}[Perfect matchings in regular bipartite graphs]
\lean{SimpleGraph.marriage_theorem}
\label{SimpleGraph.marriage_theorem}
Let $G$ be a simple graph on a finite vertex set that is bipartite with parts $X$ and
$Y$, so that every edge of $G$ joins a vertex of $X$ to a vertex of $Y$. If there is an
integer $k > 0$ such that every vertex of $G$ has degree exactly $k$, then $G$ admits a
perfect matching: a set of edges that is vertex-disjoint and covers every vertex of $G$.
\end{theorem}

\begin{theorem}[Matchings are no larger than coverings]
\lean{SimpleGraph.matching_card_le_covering_card}
\label{SimpleGraph.matching_card_le_covering_card}
\uses{SimpleGraph.IsCovering}
Let $G$ be a finite simple graph. If $M$ is a matching in $G$ and $K$ is a covering of
$G$, then the number of edges of $M$ is at most the number of vertices in $K$; that is,
$\abs{M} \le \abs{K}$.
\end{theorem}

\begin{theorem}[Equal size forces maximum matching and minimum covering]
\lean{SimpleGraph.isMaximumMatching_and_isMinimumCovering_of_card_eq}
\label{SimpleGraph.isMaximumMatching_and_isMinimumCovering_of_card_eq}
\uses{SimpleGraph.IsCovering, SimpleGraph.IsMinimumCovering, SimpleGraph.Subgraph.IsMaximumMatching}
Let $G$ be a finite simple graph, let $M$ be a matching in $G$, and let $K$ be a
covering of $G$. If $|M| = |K|$, then $M$ is a maximum matching and $K$ is a minimum
covering of $G$.
\end{theorem}

\begin{theorem}[König's theorem: maximum matching equals minimum covering]
\lean{SimpleGraph.konig_matching_covering}
\label{SimpleGraph.konig_matching_covering}
\uses{SimpleGraph.IsMinimumCovering, SimpleGraph.Subgraph.IsMaximumMatching}
Let $G$ be a finite bipartite graph with bipartition $(X, Y)$ of its vertex set. If $M$
is a maximum matching of $G$ and $K$ is a minimum covering of $G$, then the number of
edges in $M$ equals the number of vertices in $K$.
\end{theorem}

\begin{theorem}[Tutte's theorem on perfect matchings]
\lean{SimpleGraph.tutte_perfect_matching}
\label{SimpleGraph.tutte_perfect_matching}
Let $G$ be a simple graph on a finite vertex set $V$. For a set $S \subseteq V$, write
$G - S$ for the graph obtained by deleting the vertices of $S$ (together with their
incident edges), and let $o(G - S)$ denote the number of odd components of $G - S$, that
is, the connected components with an odd number of vertices. Then $G$ has a perfect
matching if and only if
\[o(G - S) \le \abs{S} \qquad \text{for every } S \subseteq V.\]
\end{theorem}

\begin{theorem}[Petersen's theorem on perfect matchings of bridgeless cubic graphs]
\lean{SimpleGraph.petersen_three_regular_bridgeless}
\label{SimpleGraph.petersen_three_regular_bridgeless}
Let $G$ be a finite simple graph that is $3$-regular (every vertex has exactly three
neighbours) and bridgeless (no edge of $G$ is a bridge, i.e. removing any single edge
leaves the graph with the same number of connected components). Then $G$ has a perfect
matching: there is a set of edges, no two sharing a vertex, that together cover every
vertex of $G$.
\end{theorem}

\begin{definition}[Feasible vertex labelling: $l x + l y \ge w s(x,y)$]
\lean{SimpleGraph.IsFeasibleLabelling}
\label{SimpleGraph.IsFeasibleLabelling}
\leanok
Feasible vertex labelling: \texttt{l x + l y $\ge$ w s(x,y)}.

We define a \emph{feasible vertex labelling} as a real-valued function $l$ on the
vertex set $X \cup Y$ such that, for all $x \in X$ and $y \in Y$
\[l(x) + l(y) \geq w(xy) \qquad (5.11)\]
(The real number $l(v)$ is called the \emph{label} of the vertex $v$.) [...] No
matter what the edge weights are, there always exists a feasible vertex
labelling; one such is the function $l$ given by
\[l(x) = \max_{y \in Y} w(xy) \ \text{if} \ x \in X, \qquad l(y) = 0 \ \text{if} \ y \in Y \qquad (5.12)\]
\end{definition}

\begin{definition}[The equality subgraph $G_l$]
\lean{SimpleGraph.equalitySubgraph}
\label{SimpleGraph.equalitySubgraph}
The \textbf{equality subgraph} \texttt{G\_l}.

If $l$ is a feasible vertex labelling, we denote by $E_l$ the set of those edges
for which equality holds in (5.11); that is
\[E_l = \{xy \in E \mid l(x) + l(y) = w(xy)\}\]
The spanning subgraph of $G$ with edge set $E_l$ is referred to as the \emph{equality
subgraph} corresponding to the feasible vertex labelling $l$, and is denoted by
$G_l$.
\end{definition}

\begin{definition}[Weight of a matching]
\lean{SimpleGraph.matchingWeight}
\label{SimpleGraph.matchingWeight}
\leanok
Weight of a matching.

The book uses the notation \texttt{w(M)} without a display, in (5.13)--(5.14):

\[w(M^*) = \sum_{e \in M^*} w(e) = \sum_{v \in V} l(v) \qquad (5.13)\]
\end{definition}

\begin{definition}[An optimal matching: a max-weight perfect matching]
\lean{SimpleGraph.IsOptimalMatching}
\label{SimpleGraph.IsOptimalMatching}
\leanok
\uses{SimpleGraph.matchingWeight}
An \textbf{optimal matching}: a max-weight perfect matching.

The optimal assignment problem is clearly equivalent to that of finding a maximum-weight
perfect matching in this weighted graph. We shall refer to such a matching as an
\emph{optimal matching.}
\end{definition}

\begin{theorem}[A perfect matching in the equality subgraph is optimal]
\lean{SimpleGraph.optimal_of_perfectMatching_equalitySubgraph}
\label{SimpleGraph.optimal_of_perfectMatching_equalitySubgraph}
\uses{SimpleGraph.IsFeasibleLabelling, SimpleGraph.IsOptimalMatching, SimpleGraph.equalitySubgraph}
Let $G$ be a bipartite graph on a finite vertex set with bipartition into parts $X$ and
$Y$ whose union is the whole vertex set, and let $w$ assign a real weight to each edge.
Suppose $l$ is a feasible vertex labelling for these weights, and let $G_l$ be the
corresponding equality subgraph. If $G_l$ contains a perfect matching $M^*$, then $M^*$,
regarded as a matching of $G$, is an optimal matching of $G$ — that is, a perfect
matching of maximum total weight.
\end{theorem}

\begin{theorem}[The $k$-cube has a perfect matching]
\lean{SimpleGraph.cube_isPerfectMatching}
\label{SimpleGraph.cube_isPerfectMatching}
\uses{SimpleGraph.cube}
For every integer $k \ge 2$, the $k$-cube $Q_k$ admits a perfect matching: there is a
set of edges of $Q_k$, no two sharing a vertex, that together cover every vertex.
\end{theorem}

\begin{theorem}[Number of perfect matchings of the complete graph]
\lean{SimpleGraph.card_perfectMatching_completeGraph}
\label{SimpleGraph.card_perfectMatching_completeGraph}
For every natural number $n$, the number of perfect matchings of the complete graph
$K_{2n}$ on $2n$ vertices — that is, the number of ways to partition its vertex set into
$n$ disjoint pairs, each joined by an edge — is
\[
\frac{(2n)!}{2^{n}\,n!}.
\]
\end{theorem}

\begin{theorem}[Perfect matchings of $K_{n,n}$]
\lean{SimpleGraph.card_perfectMatching_completeBipartite}
\label{SimpleGraph.card_perfectMatching_completeBipartite}
For every natural number $n$, the complete bipartite graph $K_{n,n}$, whose two vertex
classes each have $n$ vertices and whose edges are exactly the pairs joining a vertex of
one class to a vertex of the other, has exactly $n!$ perfect matchings.
\end{theorem}

\begin{theorem}[A tree has at most one perfect matching]
\lean{SimpleGraph.isTree_subsingleton_perfectMatching}
\label{SimpleGraph.isTree_subsingleton_perfectMatching}
Let $G$ be a tree on a finite vertex set — a connected, acyclic simple graph. Then any
two perfect matchings of $G$ coincide; that is, $G$ has at most one perfect matching.
\end{theorem}

\begin{theorem}[The complete bipartite graph $K_{n,n}$ is 1-factorable]
\lean{SimpleGraph.completeBipartite_one_factorable}
\label{SimpleGraph.completeBipartite_one_factorable}
\uses{SimpleGraph.IsKFactorable}
For every natural number $n$, the complete bipartite graph $K_{n,n}$, whose two vertex
classes each have $n$ vertices and whose edges join every vertex of one class to every
vertex of the other, is $1$-factorable: its edge set decomposes into edge-disjoint
$1$-factors (perfect matchings).
\end{theorem}

\begin{theorem}[The complete graph on an even number of vertices is 1-factorable]
\lean{SimpleGraph.completeGraph_even_one_factorable}
\label{SimpleGraph.completeGraph_even_one_factorable}
\uses{SimpleGraph.IsKFactorable}
For every natural number $n$, the complete graph $K_{2n}$ on $2n$ vertices is
$1$-factorable; that is, its edge set can be partitioned into edge-disjoint $1$-factors.
\end{theorem}

\begin{theorem}[A 3-factor from a minimum-degree bound]
\lean{SimpleGraph.exists_three_factor_of_minDegree}
\label{SimpleGraph.exists_three_factor_of_minDegree}
\uses{SimpleGraph.Subgraph.IsKFactor}
Let $G$ be a finite simple graph on $n$ vertices, and suppose that $n$ is even and that
the minimum degree $\delta$ of $G$ satisfies $n + 2 \le 2\delta$, that is, $\delta \ge
\tfrac{n}{2} + 1$. Then $G$ has a $3$-factor: a spanning subgraph in which every vertex
has degree exactly $3$.
\end{theorem}

\begin{theorem}[Walecki's decomposition of $K_{2n+1}$ into connected 2-factors]
\lean{SimpleGraph.completeGraph_odd_two_factorable}
\label{SimpleGraph.completeGraph_odd_two_factorable}
\uses{SimpleGraph.Subgraph.IsKFactor}
For every integer $n \ge 1$, the complete graph $K_{2n+1}$ on $2n+1$ vertices can be
expressed as the edge-disjoint union of $n$ connected $2$-factors. Concretely, there
exist spanning subgraphs $H_1, \ldots, H_n$ of $K_{2n+1}$, each of which is $2$-regular
(every vertex has exactly two neighbours in it) and connected, whose edge sets are
pairwise disjoint and whose union is all of $K_{2n+1}$.
\end{theorem}

\begin{theorem}[Hall-type criterion for a perfect matching in a bipartite graph]
\lean{SimpleGraph.bipartite_perfectMatching_iff}
\label{SimpleGraph.bipartite_perfectMatching_iff}
Let $G$ be a simple graph on a finite vertex set $V$ that is bipartite with parts $X$
and $Y$: the sets $X$ and $Y$ are disjoint, together they exhaust $V$ (so $X \cup Y =
V$), and every edge of $G$ joins a vertex of $X$ to a vertex of $Y$. For a set $S
\subseteq V$ write $N(S) = \bigcup_{v \in S} N(v)$ for the set of all vertices adjacent
to some vertex of $S$. Then $G$ has a perfect matching if and only if $\abs{S} \le
\abs{N(S)}$ for every set $S \subseteq V$.
\end{theorem}

\begin{theorem}[Every regular bipartite graph is 1-factorable]
\lean{SimpleGraph.regular_bipartite_one_factorable}
\label{SimpleGraph.regular_bipartite_one_factorable}
\uses{SimpleGraph.IsKFactorable}
Let $G$ be a bipartite graph with vertex classes $X$ and $Y$, and let $k$ be a positive
integer. If $G$ is $k$-regular, then $G$ is $1$-factorable; that is, $G$ decomposes into
edge-disjoint $1$-factors.
\end{theorem}

\begin{theorem}[Petersen's two-factorization of even-regular graphs]
\lean{SimpleGraph.regular_two_factorable}
\label{SimpleGraph.regular_two_factorable}
\uses{SimpleGraph.IsKFactorable}
Let $G$ be a finite simple graph, and suppose $G$ is regular of degree $2k$ for some
natural number $k$, meaning every vertex has exactly $2k$ neighbours. Then $G$ is
$2$-factorable: its edges decompose into finitely many pairwise edge-disjoint
$2$-factors, that is, edge-disjoint $2$-regular spanning subgraphs whose union is all of
$G$.
\end{theorem}

\begin{theorem}[König's theorem in line-cover form for 0–1 matrices]
\lean{SimpleGraph.konig_matrix}
\label{SimpleGraph.konig_matrix}
Let $A$ be an $m \times n$ matrix with entries in $\{0,1\}$, and call a *line* of $A$
any single row or single column. Consider two quantities: the minimum number $c$ of
lines needed to cover all the $1$-entries of $A$ (that is, the least value of $\abs{R} +
\abs{C}$ over all sets $R$ of rows and $C$ of columns such that every entry equal to $1$
lies in a row of $R$ or a column of $C$), and the maximum number $d$ of $1$-entries of
$A$ no two of which share a line (that is, the largest size of a set of positions with
entry $1$ whose members have pairwise distinct rows and pairwise distinct columns). Then
$c = d$.
\end{theorem}

\begin{theorem}[König–Ore defect formula for bipartite matchings]
\lean{SimpleGraph.konig_ore_defect}
\label{SimpleGraph.konig_ore_defect}
\uses{SimpleGraph.Subgraph.IsMaximumMatching}
Let $G$ be a finite bipartite graph with bipartition $(X, Y)$, and for a set $S
\subseteq X$ write $N(S) = \bigcup_{v \in S} N(v)$ for the set of vertices adjacent to
some vertex of $S$. If $M$ is a maximum matching of $G$, then
\[
\lvert M\rvert \;+\; \max_{S \subseteq X}\,\max\bigl(\lvert S\rvert - \lvert
N(S)\rvert,\;0\bigr) \;=\; \lvert X\rvert,
\]
where $\lvert M\rvert$ denotes the number of edges of $M$.
\end{theorem}

\begin{theorem}[Large edge count forces a matching of size $k$]
\lean{SimpleGraph.exists_matching_card_eq_of_card_edges_gt}
\label{SimpleGraph.exists_matching_card_eq_of_card_edges_gt}
Let $G$ be a finite bipartite graph with vertex parts $X$ and $Y$, so that every edge of
$G$ joins a vertex of $X$ to a vertex of $Y$, and suppose $\abs{X} = \abs{Y} = n$. Write
$\varepsilon$ for the number of edges of $G$, and let $k$ be a natural number. If
$(k-1)\,n < \varepsilon$, then $G$ contains a matching consisting of exactly $k$ edges.
\end{theorem}

\begin{theorem}[Doubly balanced matrices are square]
\lean{SimpleGraph.doublyStochastic_card_eq}
\label{SimpleGraph.doublyStochastic_card_eq}
Let $m$ and $n$ be finite index sets, and let $Q = (Q_{ij})_{i \in m,\, j \in n}$ be a
real matrix whose every row sums to $1$, so that $\sum_{j \in n} Q_{ij} = 1$ for each $i
\in m$, and whose every column sums to $1$, so that $\sum_{i \in m} Q_{ij} = 1$ for each
$j \in n$. Then the two index sets have the same size, $\abs{m} = \abs{n}$; that is, $Q$
is square.
\end{theorem}

\begin{theorem}[Common transversal of left and right cosets (P. Hall)]
\lean{SimpleGraph.exists_common_coset_transversal}
\label{SimpleGraph.exists_common_coset_transversal}
Let $H$ be a finite group and let $K \le H$ be a subgroup. Then there exist a natural
number $n$ and elements $h_1, \dots, h_n \in H$ such that the assignment $i \mapsto h_i
K$, sending each index to the corresponding left coset viewed as a subset of $H$, is a
bijection from $\{1, \dots, n\}$ onto the collection of all subsets of $H$, and,
simultaneously, the assignment $i \mapsto K h_i$, sending each index to the
corresponding right coset, is also a bijection from $\{1, \dots, n\}$ onto the
collection of all subsets of $H$.
\end{theorem}

\begin{theorem}[Perfect matchings in edge-connected regular graphs]
\lean{SimpleGraph.perfectMatching_of_edgeConnected_regular}
\label{SimpleGraph.perfectMatching_of_edgeConnected_regular}
\uses{SimpleGraph.edgeConnectivity}
Let $G$ be a finite simple graph that is $k$-regular for some integer $k \ge 1$, and
whose edge connectivity satisfies $\kappa'(G) \ge k - 1$. If $G$ has an even number of
vertices, then $G$ has a perfect matching, that is, a set of pairwise disjoint edges
that together cover every vertex.
\end{theorem}

\begin{theorem}[Perfect matchings in trees via odd components]
\lean{SimpleGraph.isTree_perfectMatching_iff}
\label{SimpleGraph.isTree_perfectMatching_iff}
Let $G$ be a finite tree, and for a graph $H$ write $o(H)$ for the number of connected
components of $H$ having an odd number of vertices. Then $G$ has a perfect matching if
and only if $o(G - v) = 1$ for every vertex $v$ of $G$, where $G - v$ denotes the graph
obtained from $G$ by deleting $v$ together with all edges incident to it.
\end{theorem}

\begin{theorem}[Berge's defect formula for maximum matchings]
\lean{SimpleGraph.berge_tutte_defect}
\label{SimpleGraph.berge_tutte_defect}
\uses{SimpleGraph.Subgraph.IsMaximumMatching}
Let $G$ be a finite simple graph on vertex set $V$, and let $M$ be a maximum matching of
$G$, so that no matching of $G$ has more edges than $M$. For a set $S \subseteq V$,
write $o(G - S)$ for the number of connected components having an odd number of vertices
in the subgraph of $G$ induced on $V \setminus S$. Then
\[
2\,\abs{M} + \max_{S \subseteq V}\bigl(o(G - S) - \abs{S}\bigr) = \abs{V},
\]
where $\abs{M}$ is the number of edges of $M$ and the maximum ranges over all subsets
$S$ of $V$.
\end{theorem}

\begin{theorem}[Large edge count forces a perfect matching]
\lean{SimpleGraph.perfectMatching_of_card_edges_gt}
\label{SimpleGraph.perfectMatching_of_card_edges_gt}
Let $G$ be a finite simple graph on $n$ vertices with $m$ edges and minimum degree
$\delta$. Suppose that $n$ is even, that $2\delta < n$, and that
\[ m > \binom{\delta}{2} + \binom{n - 2\delta - 1}{2} + \delta\,(n - \delta). \]
Then $G$ has a perfect matching.
\end{theorem}
